% Shared look for every deck: Metropolis with a progress bar under each frame title.
% Build with Tectonic: `make` in this directory.

\usetheme[progressbar=frametitle,numbering=fraction,sectionpage=progressbar,block=fill]{metropolis}

% Metropolis looks Fira Sans up by system name; load it from the TeX bundle instead,
% Light for body text as Metropolis intends, with lining figures for data.
\usepackage[sfdefault,light,lf]{FiraSans}
\usepackage{booktabs,tabularx}
\usepackage{tikz}
\usetikzlibrary{arrows.meta,positioning,shapes.geometric,calc}

% Metropolis's palette, with the decks' own colour names mapped onto it.
\definecolor{mDarkTeal}{HTML}{23373B}
\definecolor{mDarkBrown}{HTML}{604C38}
\definecolor{mLightBrown}{HTML}{EB811B}
\colorlet{crust}{mDarkBrown}
\definecolor{crumb}{HTML}{FAEEE2}
\colorlet{rye}{mDarkTeal}
\definecolor{flag}{HTML}{B03A2E}
\definecolor{own}{HTML}{2E7D4F}
\definecolor{muted}{HTML}{7A7F80}

% Top-align frames, so the \vfill in \source puts the source line at the bottom.
% A slimmer title bar and narrower margins than Metropolis's defaults leave room
% for the denser slides.
\makeatletter
\beamer@centeredfalse
\setlength{\metropolis@frametitle@padding}{1.4ex}
% A progress bar you can see: under every frame title and on section pages.
\setlength{\metropolis@progressinheadfoot@linewidth}{2pt}
\setlength{\metropolis@progressonsectionpage@linewidth}{1.5pt}
\setlength{\metropolis@titleseparator@linewidth}{1pt}
% Metropolis's title page is \paperheight tall, which overflows a top-aligned frame.
\setbeamertemplate{title page}{%
  \begin{minipage}[b][\textheight]{\textwidth}
    \vfill
    \usebeamertemplate*{title}
    \ifx\insertsubtitle\@empty\else\usebeamertemplate*{subtitle}\fi
    \usebeamertemplate*{title separator}
    \usebeamertemplate*{author}
    \usebeamertemplate*{institute}
    \vfill
  \end{minipage}}
\makeatother
\setbeamersize{text margin left=1.2em,text margin right=1.2em}

\setbeamertemplate{frame footer}{\insertshorttitle}
\setbeamertemplate{footline}{%
  \begin{beamercolorbox}[wd=\paperwidth,ht=2.6ex,dp=1.2ex,leftskip=1.2em,rightskip=1.2em]{footline}%
    \usebeamerfont{page number in head/foot}%
    \usebeamertemplate*{frame footer}\hfill\usebeamertemplate*{frame numbering}%
  \end{beamercolorbox}}
\setbeamerfont{page number in head/foot}{size=\tiny}
\setbeamercolor{footline}{fg=muted}

% The practical point of a slide, for the baker.
\newcommand{\takeaway}[1]{%
  \par\medskip
  \colorbox{crumb}{\parbox{\dimexpr\linewidth-2\fboxsep}{\small\textbf{\color{crust}In the kitchen:} #1}}}

% Small source line, placed at the bottom of a frame.
\newcommand{\source}[1]{%
  \vfill{\color{muted}\tiny Source: #1}}

% Status flags, mirroring the notes' conventions. Filled tags, so they read on the
% dark title bar as well as on the slide body.
\newcommand{\flagtag}[2]{\mbox{\setlength{\fboxsep}{1.2pt}\colorbox{#1}{\color{white}\scriptsize\bfseries\,#2\,}}}
\newcommand{\hypothesis}{\flagtag{flag}{hypothesis}}
\newcommand{\contested}{\flagtag{flag}{contested}}
\newcommand{\ownresult}{\flagtag{own}{own result}}
\newcommand{\untested}{\flagtag{flag}{not yet tested}}

% Source lists: one source per bullet, author and year in bold, optional group
% headings. Use in a [allowframebreaks] frame; long lists continue on "(cont.)" slides.
\newenvironment{sourcelist}{%
  \setbeamerfont{itemize/enumerate body}{size=\scriptsize}%
  \setbeamerfont{itemize item}{size=\scriptsize}%
  \scriptsize\begin{itemize}\setlength{\itemsep}{2pt}}{\end{itemize}}
\newcommand{\src}[2]{\item\textbf{#1}\enspace #2}
% \nobreak keeps a group heading on the same slide as its first source.
\newcommand{\srcgroup}[1]{\par\smallskip{\small\bfseries\color{crust}#1}\par\nobreak\vspace{-2pt}\nobreak}
\setbeamertemplate{frametitle continuation}[from second][(cont.)]

% Flowchart nodes for TikZ diagrams.
\tikzset{
  box/.style={draw=crust,fill=crumb,rounded corners=2pt,align=center,
    font=\scriptsize,inner sep=4pt,text width=#1},
  box/.default=2.6cm,
  arrow/.style={-{Stealth[length=5pt]},draw=crust,thick},
}

\author{Sourdough Science}
\institute{github.com/zrrbite/sourdough-science}
\date{}
